\documentclass[11pt,a4paper]{article}
\usepackage{fontspec,amsmath,amssymb,amsthm,microtype}
\setmainfont{lmroman10-regular.otf}[BoldFont=lmroman10-bold.otf,ItalicFont=lmroman10-italic.otf,BoldItalicFont=lmroman10-bolditalic.otf]
\setmonofont{lmmono10-regular.otf}[Scale=MatchLowercase]
\usepackage[margin=25mm]{geometry}
\usepackage[unicode,hidelinks]{hyperref}
\setlength{\emergencystretch}{3em}
\setlength{\parskip}{0.4em}
\setlength{\parindent}{0pt}
\allowdisplaybreaks[2]
\raggedbottom
\title{\texorpdfstring{Kernel-\allowbreak{}checked declaration and axiom audit}{Kernel-checked declaration and axiom audit}}
\author{LangevinLean mathematical exposition}
\date{}
\begin{document}\maketitle
\noindent Source module: \texttt{Proximal\allowbreak{}Langevin/\allowbreak{}Check.\allowbreak{}lean}. The accompanying HTML pairs this mathematical exposition with the exact Lean statements and records the scope of the formalization.


This module imports the proximal development and asks Lean to print the axioms used by its principal analytic, algorithmic and stochastic declarations. It is an audit entry point rather than a new mathematical assumption.


\section{\texorpdfstring{Rebuild the theorems and inspect their dependency axioms}{Rebuild the theorems and inspect their dependency axioms}}


Run \texttt{lake build Proximal\allowbreak{}Langevin +Proximal\allowbreak{}Langevin.\allowbreak{}Check} from the repository root. Each \texttt{\#print axioms} command reports the dependency closure of an already checked declaration, including proximal optimality, the three actual algorithm bounds, Brownian coupling and mean-field constructions. The file introduces no new theorem or definition. The audit complements compilation; it does not strengthen any theorem’s stated assumptions or turn conditional scalar statements into algorithm convergence results.



Each \texttt{\#check} command asks Lean to elaborate the complete theorem type, retaining every displayed assumption. Each \texttt{\#print axioms} command reports the transitive foundational dependencies of that theorem. Compile this module after the library build to inspect that report; these commands check the existing proofs and do not supply new mathematical hypotheses or conclusions.


\end{document}
